\section*{ID9204: Definition of P)), BerryPhase}
\paragraph{Metadata.}
PiID: 2103274889218 | RTEID: RTE:P35741.1760:T2.3418 | UUID: 7b261953-8e90-5ae2-8994-eb18cf74aaf1

\paragraph{Canonical Statement.}
Handle: Perform experiments isolating each stage. For example, apply the rotating dipoles without layering and verify that thrust does not occur; such tests can falsify individual stage requirements. 2.0.70 ID 219: QuantumCoriolis Operator Type: Operator Formal Statement: The quantumCoriolis operator for a state psi with angular momentum operator J and rotation vector Omega is defined as the triple \{UnitaryEvolution = exp(-i H\_rot t/hbar), DensityBifurcation = TensorDecompose(-2 Omega × (rho \$⊗\$ p)), BerryPhase = integral from 0 to 2 pi/Omega of 〈psi(t)| i d\_t |psi(t)〉 dt\}. Here H\_rot = -hbar Omega·J and rho = psi \$⊗\$ psi. Interpretive Meaning: This operator bundles together the rotational Hamiltonian, the Coriolisinduced density matrix bifurcation and the Berry phase accumulated during rotation. It provides a unifying framework for analysing quantum systems under rotation and Coriolis forces. Canonical Registry Vol. 01 74 Trawin Zero Point Contextus Canonicus 2 TZP CANONICAL REGISTRY

\paragraph{Canonical Equation.}
\[
p)), BerryPhase = integral from 0 to 2 pi/Omega of 〈psi(t)| i d_t |psi(t)〉 dt}. Here H_rot = -hbar Omega·J and rho = psi
\]

\paragraph{Dictionary.}
Handle: Perform experiments isolating each stage.

\paragraph{Encyclopedia.}
Definition of P)), BerryPhase is an algebraic definition, written p)), BerryPhase = integral from 0 to 2 pi/Omega of 〈psi(t)| i d\_t |psi(t)〉 dt. H. It is an algebraic definition expressing the quantity directly in terms of the variables on its right-hand side. It is built from n, g, r, a, l, f, defining p)), BerryPhase. Within the Trawin Zero Point registry it serves as one element of the framework's mathematical narrative.
